% Tapis de Quillen, with Quillen's offprint and preprint (cote 162-5, typed part): archive cote 162-5, site page 8 (PDF page 9).
% typescript. Draft compared: none.
% AI-transcribed by Claude Opus 5.5 (claude-opus-5-5, adaptive thinking, effort high), Message Batches API, batch pilotnext (msgbatch_01A9sdu8SYq23bEtqX9tUncQ), 2026-11-13; not yet checked by eye.
\scanpage{162-5/8}
\margin{\unsure{!} Attention à la catégorie \uline{Hom} \hand{\unsure{7}} grad., qui s'identififie à la catégorie \unsure{définie} \unsure{par un} $R\mathrm{Hom}_{\mathbb{Z}}(C_{\bullet}, C'_{\bullet})$, ce point devrait être élucidé!}

de même les (Cat)-groupes(ou groupes à valeurs dans (Cat)) comme étant
les groupes semisimpliciaux qui "sont"des catégories i.c. qui \xout{sexdéfinid}
\xout{sent} sat\unsure{is}font la condition d'exactitude familière, permettant d'en déte
miner tous les composants à l'aide de $C_{o\xout{m}}$ et $C_1$ . Appliquant le tapis de
Dold, dans le cas commutatif, on trouve que les (Cat)-groupes abéliens
forment une catégorie équivalente à celle des complexes de chaines de
groupes abéliens qui sont de longueur 1:
\[ 0 \leftarrow C_o \leftarrow C_1 \leftarrow 0 \quad . \]
Quand on explicite la (Cat)-groupe abélien associé à un tel complexe de
chaines, on trouve "ma" construction chérie da la "catégorie de Picard"
associée au complexe: objets ceux de $C_o$ , flêches de x dans y les z dans
tels que $dz = y-x$ . On trouve donc à isomorphisme \unsure{près|ppès} tous les group
abéliens de (Cat) ainsi. Un complément important est que toute catégori
\margin{\unsure{Objections déjà} \unsure{mais cela contredit} \unsure{ce qui est dit des} \illegible{2} \unsure{Picard} \illegible{3} \unsure{Mais} \illegible{2} \unsure{de Picard} \illegible{4} \unsure{(Cat)} \illegible{2}}
de Picard \inserted[typed]{(resp. \deleted{et}} commu\deleted{tative)}est équivalente (au sens technique qu'on ne précis
pas ici) à \xout{laxcatégoxxie} un vrai (Cat)-groupe \inserted[typed]{\deleted{(resp. (Cat)-groupe}} \deleted{(commutatif)} . \hand{Quand on}
\hand{\unsure{identifie} mod ``\unsure{Gr}-équivalence'', \unsure{cela} correspond à travailler dans la catégorie \uline{dérivée}}
\hand{\unsure{des complexes}, \unsure{seuls invariants} \unsure{étant} $\pi_0, \pi_1$)}
On notera qu'on a une \xout{lxxxxx} classification beaucoup plus restrein
que pour les H-espaces (resp. les H-espaces homotopiquement comm.) mod
équivalence d'homotopie, de façon précise entre les groupes resp. les gr
oupes abéliens de la catégorie homotopique. La raison en est quon est beau
coup plus restrictif dans le contexte catégorique pour le sens da l'asso
tivité (resp. et de la commutativité), puisqu'on ne demande pas seulemen
l'existence d'isomorphismes d'associ\unsure{ai}tivité\unsure{$_x$} (qui suffirait pour avoir u
groupe dans \inserted[typed]{la catégorie homotopique} (Cat)' $\simeq$ (Ssimpl)', modulo l'existence d'unités et d'invers
mais qu'on suppose \uline{données} ces isomorphismes, de façon à satisfaire à ce
taines compatibilités: c'est ce qui permet, me semble-t-il, de prouver qu
\margin{(}équivalence près, on peut supposer qu'il s'agit de vrais groupes dans
(Cat), \deleted{resp. de vraie commutativité.}

\mspage{8}
